% !TeX root = ../概率论与数理统计（缪柏其）笔记.tex
% 由主文档通过 \input 载入；请编译主文档。
\chapter{绪论}

概率论与统计学都研究随机现象和不确定性，但侧重点不同。\emph{概率论（probability theory）}从给定的假设和规则出发，研究各种结果出现的可能性及其规律；\emph{统计学（statistics）}研究如何收集和利用数据，描述已经观察到的现象、判断未知情况，并评价判断的可靠程度。\emph{前者侧重从模型推导可能的结果，后者的一个主要任务是利用数据认识模型及其未知部分。}

以重复抛掷硬币为例，假定每次正面出现的机会保持不变，各次抛掷的结果互不影响。如果还知道正反面出现的机会相同，研究“抛掷 100 次，恰好出现 60 次正面的可能性有多大”，就是典型的概率论问题。如果正面出现的机会未知，实际抛掷 100 次，记录到 60 次正面，再据此估计这一机会，并考察估计可能有多大误差，就是典型的统计学问题。

\emph{“记录中正面占六成”是在描述已观察到的事实，并不意味着硬币每次正面出现的真实机会一定是六成。}即使正反面的机会相同，抛掷 100 次也不保证正面恰好出现 50 次。因此，根据数据判断硬币的性质时，需要考虑偶然波动，不能仅凭次数偏离 50 次就断言硬币不公平。

概率论为统计方法提供数学基础，使我们能够研究数据和判断中的随机波动；统计学则利用数据认识未知情况，并检查模型是否适合实际问题。\emph{数理统计（mathematical statistics）}着重研究统计方法的数学原理及其性质，是统计学的重要组成部分。两者合在一起学习，体现了从认识随机规律到利用数据作出判断的知识联系。\emph{学习时应分清模型的假设、实际的观察记录，以及据此能够得到的结论。}

本笔记以缪柏其、张伟平主编的《概率论与数理统计》（高等教育出版社，2022 年第 1 版）为学习线索，按教材章节顺序整理学习思考与习题解答。



